% =============================================================
% INTRODUCTION (unnumbered)
% =============================================================
\section*{Introduction}

On the eve of the First World War a visitor to Riga or Tallinn found a `culturally Western' province of the Russian Empire, where people worshipped in Lutheran churches, read in their own languages, traded with the West and enjoyed incomes closer to Europe's than to Russia's; the imperial government called the Baltic provinces `our little West' \parencites[129]{corum2013}[427]{kodres_soviet_west}. In 1913 income per head in Latvia and Estonia stood at the level of Finland or Italy, and even agrarian Lithuania was richer than Greece or Portugal. When Red Army soldiers entered the Baltic republics under the mutual-assistance pacts of 1939 they wrote home in astonishment: `the store is full of all sorts of fabrics, whatever the soul desires, and there is not a single person, not to mention a line' \parencite[279]{meltyukhov2014en}. By 1990, after forty-five years inside the Soviet Union, income per head in the three republics was about half that of Finland and three-fifths that of the Western periphery; by 2024, after three decades of reform, European Union membership and `Baltic tiger' growth, they remained in Europe's periphery, and the richest of them in 1913, Latvia, had become the poorest. Three small, literate, egalitarian economies, which possessed from the outset much of what the catch-up literature regards as the social capability needed to grow, have failed over 110 years and under three economic regimes to leave the periphery into which the nineteenth century placed them. That is the Baltic growth puzzle.

The puzzle has not been posed in these terms because the data to pose it did not exist. The Maddison Project Database, the standard source for long-run comparisons, reports GDP per capita for Lithuania, Latvia and Estonia only for 1973 and from 1980, and its pre-1990 figures descend from Soviet net material product, which Maddison himself regarded as unusable \parencites[307--8]{Maddison1998b}{mpd2023}. The Baltic states have accordingly been absent from the economic histories of twentieth-century Europe, and discussion of their living standards has been left to memory and politics: a romanticised interwar `golden age' on one side, and on the other the claim, still made, that Soviet rule brought no progress at all \parencites[92]{lehti_new_lvee}[2]{ehbaltics}. The gap has begun to close. Direct benchmark estimates now exist for Lithuania in 1937, Latvia in 1935 and Estonia in 1938, annual interwar series have been built on them, and the first Soviet-period series, for Lithuania and Estonia, have shown that the Maddison Project's late-Soviet levels are too high by between a sixth and a third \parencite{klimantas2019,norkuslatvia1935,Klimantas2024,klimantas2026soviet,klimantas2025estonia}. Latvia is the missing piece. Its interwar record is the most disputed of the three, and for the Soviet period there is no estimate of real output at all: the Western Sovietologists who studied it worked from official aggregates, and Latvian historians, who have reconstructed the institutional history of the Soviet economy in detail, have had to rely on indices whose growth is several times the real one \parencite{king_1965_economic,viksnins_1981_evaluating,gatis_lv_eh}.

This paper makes three contributions. First, it provides the first annual series of Latvian GDP for the Soviet period, 1937--95, built from 129 physical product series, agricultural output, transport volumes and employment records drawn from Soviet yearbooks, the German occupation statistics of 1942--4 and the Latvian and Russian archives, anchored to Bolotin's purchasing-power-parity benchmark for 1990 with the methods of the CIA's accounts for the Soviet Union and of the companion Lithuanian and Estonian series. Second, it splices the new series to the published interwar series, extended here to 1919 and 1940, and to adjusted post-1991 data, to give a continuous series for 1913 and 1919--2024 in 1990 and 2011 international dollars, with a consistent population series, ready for the Maddison Project Database. Because the interwar and Soviet-period series are built from different sources and anchored to benchmarks fifty-five years apart, their agreement where they overlap, within 3.5 per cent in each of the years 1937--40 and within 3 per cent at the 1990--1 splice, is the strongest available check on the whole reconstruction. Third, and principally, it uses the completed Baltic record to provide the first comparison of the three Baltic economies, with each other and with the rest of Europe, that rests on national accounts rather than proxies, across interwar capitalism, Soviet socialism and restored independence. The Latvian series is analysed briefly, since the companion papers have examined the Lithuanian and Estonian cases in depth; the weight of the paper falls on the comparison.

The new series alters the received picture in four respects. First, the Baltic states did not converge on Western Europe under interwar capitalism. Latvia and Lithuania emerged from the First World War as the most devastated economies in Europe, having lost, like Russia, three-fifths of their output, with income per head more than a third below 1913; the 1920s were spent regaining it, and the rapid growth of 1929--38, \gr{1929}{1938} per cent a year in Latvia and almost as much in its neighbours, faster than anywhere in Europe outside Finland and the Soviet Union, only restored the region to \pctx{Baltics}{Western}{1938} per cent of the Western core, where it had stood in 1913. Second, the Soviet period divides in two. Between 1950 and 1968 income per head grew by \gr{1950}{1968} per cent a year in Latvia and faster in Lithuania, and the region reached \pctx{Baltics}{Western}{1968} per cent of the Western core, closer than at any time before 2019; between 1968 and 1989 Latvia grew by \gr{1968}{1989} per cent a year, slower than any economy of the communist bloc for which a series exists, and the gains were lost. Within the Union the Baltic republics remained the richest throughout, Latvia at \pctussr{1968} per cent of the Soviet average in 1968 and \pctussr{1989} in 1989; the lead was inherited rather than created by Soviet investment, which as a share of the republic's GDP stayed below the Union average, and it narrowed in the final decades. Third, the transition cut Latvian income per head by \lvTROUGHpct{} per cent in four years, the deepest of the three Baltic contractions, and the convergence that followed, at \gr{1995}{2007} per cent a year to 2007, is the only sustained convergence the region has known; since 2007 it has been uneven, and Lithuania, the poorest of the three for most of the century, overtook Latvia durably in \ltOVERTAKElv{} and Estonia in \ltOVERTAKEee{}. Fourth, the belief that the Baltic states stood level with Finland before the war, which underlies most estimates of the cost of occupation, is not borne out: Latvia had fallen from \pctfin{1913} per cent of the Finnish level in 1913 to \pctfin{1938} per cent by 1938, before Soviet rule began.

The paper also speaks to a wider literature. The economic history of the Soviet Union has been rewritten from archival sources, but that work has concentrated on Russia, and the regional development of the Union, which Kukić has called a major gap in its spatial economic analysis, remains almost without national accounts \parencite{zhuravskaya2024,kukic_soviet}. The Baltic states, with borders stable since 1945 and a capitalist past against which the socialist record can be measured, offer the first regional evidence, and the procedure used here can be applied to Ukraine, Belarus and the other republics. They also add three cases to the debate on peripheral convergence in Europe, in which Eastern Europe's failure to catch up under both interwar capitalism and socialism has been attributed to deficient social capability inherited from a semi-feudal past \parencite{morysinterwarinmorys2021,vonyoinmorys2021,berendregions2020}. The Baltic states had shed that past by 1925; their failure to converge before 1990 therefore points to the regimes and the external environment rather than to the societies.

The next section reviews the data gap and the earlier estimates. The following section describes the sources and methods of the Latvian Soviet-period series, its splicing to the interwar and modern series, and the limits of the result. A brief section then reads the Latvian series. The core of the paper compares the three Baltic economies under the three regimes, with growth rates, relative levels, structure, convergence tests and a discussion of the sources of catch-up and decline. A final section concludes.
